% Part: computability
% Chapter: recursive-functions
% Section: pr-relations

\documentclass[../../../include/open-logic-section]{subfiles}

\begin{document}

\olfileid{cmp}{rec}{prr}
\olsection{بنیادی بازگشتی تعلق}


\begin{defn}
تعلق $R(\vec x)$ نوں بنیادی بازگشتی
اوہدوں کہندے نیں جدوں اوہدا خصوصیتی
فنکشن
\[
\Char{R}(\vec x) = \left\{
  \begin{array}{ll}
  1 & \mbox{جے $R(\vec x)$ درست ہووے} \\
  0 & \mbox{ورنہ}
  \end{array}
\right.
\]
بنیادی بازگشتی ہووے۔
\end{defn}

دوجے لفظاں وچ، جدوں اسی بنیادی
بازگشتی تعلق $R(\vec x)$ دی گل
کردے آں، تاں مراد اوہ تعلق اے
جیہڑا $\Char{R}(\vec x) = 1$
دے روپ وچ لکھیا جا سکدا اے،
جتھے $\Char{R}$ اک بنیادی
بازگشتی فنکشن اے جیہڑا ہر انپٹ
اُتے یا 1 یا 0 واپس دیندا اے۔
مثال لئی، تعلق $\fn{IsZero}(x)$
اوہدوں تے صرف اوہدوں درست اے
جدوں $x = 0$ ہووے؛ ایہ فنکشن
$\Char{\fn{IsZero}}$ دے مطابق اے،
جیہڑا بنیادی بازگشت راہیں
ایویں متعین اے:
\begin{align*}
\Char{\fn{IsZero}}(0) & = 1,\\
\Char{\fn{IsZero}}(x+1) & = 0.
\end{align*}

ایہ صاف ہونا چاہیدا اے کہ تعلقاں
نوں ہور بنیادی بازگشتی فنکشناں
نال ترکیب دے روپ وچ ورتیا جا
سکدا اے۔ ایس لئی ہیٹھلے تعلق
وی بنیادی بازگشتی نیں:
\begin{enumerate}
\item برابری دا تعلق، $x = y$،
  جیہڑا $\fn{IsZero}(\left|x -
  y\right|)$ راہیں متعین اے۔
\item چھوٹا یا برابر ہوون دا تعلق،
  $x \leq y$، جیہڑا $\fn{IsZero}(x
  \tsub y)$ راہیں متعین اے۔
\end{enumerate}


\begin{prop}
  بنیادی بازگشتی تعلقاں دا سیٹ
  بولیائی عملاں تحت بند اے؛ یعنی
  جے $P(\vec x)$ تے $Q(\vec x)$
  بنیادی بازگشتی نیں، تاں
  ایہ وی نیں:
  \begin{enumerate}
  \item $\lnot P(\vec x)$
  \item $P(\vec x) \land Q(\vec x)$
  \item $P(\vec x) \lor Q(\vec x)$
  \item $P(\vec x) \lif Q(\vec x)$
  \end{enumerate}
\end{prop}

\begin{proof}
  فرض کرو $P(\vec x)$ تے $Q(\vec x)$
  بنیادی بازگشتی نیں، یعنی اوہناں دے
  خصوصیتی فنکشن $\Char{P}$ تے
  $\Char{Q}$ ایہو جیہے نیں۔ سانوں
  دکھانا اے کہ $\lnot P(\vec x)$
  وغیرہ دے خصوصیتی فنکشن وی
  بنیادی بازگشتی نیں:
  \[
  \Char{\lnot P}(\vec x) = \begin{cases}
    0 & \text{جے $\Char{P}(\vec x) = 1$ ہووے}\\
    1 & \text{ورنہ}
  \end{cases}
  \]
  اسی $\Char{\lnot P}(\vec x)$
  نوں $1 \tsub \Char{P}(\vec x)$
  دے روپ وچ متعین کر سکدے آں:
  \[
  \Char{P \land Q}(\vec x) = \begin{cases}
    1 & \text{جے $\Char{P}(\vec x) = \Char{Q}(\vec x) = 1$ ہووے}\\
    0 & \text{ورنہ}
  \end{cases}
  \]
  اسی $\Char{P \land Q}(\vec x)$
  نوں $\Char{P}(\vec x) \cdot
  \Char{Q}(\vec x)$ یا
  $\fn{min}(\Char{P}(\vec x), \Char{Q}(\vec
  x))$ دے روپ وچ متعین کر سکدے آں۔
  اِس طرح،
  \begin{align*}
    \Char{P \lor Q}(\vec x) & = \fn{max}(\Char{P}(\vec x), \Char{Q}(\vec x)) \text{ تے}\\
    \Char{P \lif Q}(\vec x) & = \fn{max}(1 \tsub
  \Char{P}(\vec x), \Char{Q}(\vec x)).
  \end{align*}
\end{proof}

\begin{prop}
  بنیادی بازگشتی تعلقاں دا سیٹ
  محدود مسوّریت تحت بند اے؛ یعنی
  جے $R(\vec x, z)$ بنیادی
  بازگشتی تعلق اے، تاں ایہ تعلق
  وی بنیادی بازگشتی نیں:
  \begin{align*}
    & \bforall{z < y}{R(\vec x, z)} \text{ تے}\\
    & \bexists{z < y}{R(\vec x, z)}.
  \end{align*}
  $\bforall{z < y}{R(\vec x, z)}$
  ، $\vec x$ تے $y$ لئی اوہدوں
  تے صرف اوہدوں درست اے جدوں
  $R(\vec x, z)$ ہر اوہ~$z$
  لئی درست ہووے جیہڑا~$y$
  توں چھوٹا اے؛ اِسے طرح
  $\bexists{z < y}{R(\vec x, z)}$
  لئی وی متعلقہ وجودی شرط اے۔
\end{prop}

\begin{proof}
  رواجاً اسی $\bforall{z < 0}{R(\vec x, z)}$
  نوں درست منّدے آں (سادہ وجہ
  ایہ اے کہ کوئی $z$،~$0$
  توں چھوٹا نہیں) تے
  $\bexists{z < 0}{R(\vec x, z)}$
  نوں غلط منّدے آں۔ محدود کلی
  مسوّر محدود حاصلِ ضرب یا
  وار وار گھٹ قدر وانگ کم کردا
  اے؛ یعنی جے $P(\vec x, y) \defiff \bforall{z <
  y}{R(\vec x, z)}$، تاں
  $\Char{P}(\vec x, y)$ ایویں
  متعین کیتا جا سکدا اے:
  \begin{align*}
    \Char{P}(\vec x, 0) & = 1\\
    \Char{P}(\vec x, y+1) & =
    \fn{min}(\Char{P}(\vec x, y), \Char{R}(\vec x, y)).
  \end{align*}
  محدود وجودی مسوّریت نوں
  اِسے طرح $\fn{max}$ ورت کے
  متعین کیتا جا سکدا اے۔
  یا فیر محدود کلی مسوّریت
  توں ایہ ہم ارزی ورت کے:
  $\bexists{z < y}{R(\vec x, z)}
  \liff \lnot \bforall{z < y}{\lnot R(\vec x, z)}$۔
  غور کرو کہ، مثال لئی،
  $\bexists{x \leq y}{\dots
  x\dots}$ دے روپ دا محدود
  مسوّر $\bexists{x < y+1}{\dots x \dots}$
  دے برابر اے۔
\end{proof}

\begin{prob}
  دکھاؤ کہ تِن-جگہ تعلق $x \equiv y \mod n$
  (یعنی modulo~$n$ دے حساب نال مطابقت)
  بنیادی بازگشتی اے۔
\end{prob}

اک ہور کم آؤن والا بنیادی بازگشتی
فنکشن شرطی فنکشن $\fn{cond}(x,y,z)$
اے، جیہڑا ایویں متعین اے:
\begin{align*}
  \fn{cond}(x,y,z) & = \begin{cases}
  y & \text{جے $x = 0$ ہووے} \\
  z & \text{ورنہ}.
\end{cases}
\intertext{ایہ بازگشتی طور اُتے ایویں متعین اے}
\fn{cond}(0,y,z) & = y,\\
\fn{cond}(x+1,y,z) & = z.
\end{align*}
ایس نوں ورت کے بنیادی بازگشتی
تعلقاں راہیں صورتاں دے حساب
نال بنیادی بازگشتی فنکشناں
دیاں تعریفاں دا جواز دتا جا
سکدا اے:

\begin{prop}
جے $g_0(\vec x)$، \dots،~$g_m(\vec x)$
بنیادی بازگشتی فنکشن نیں، تے
$R_0(\vec x)$، \dots، $R_{m-1}(\vec x)$
بنیادی بازگشتی تعلق نیں، تاں
ہیٹھ تعریف کیتا فنکشن $f$
\[
f(\vec x) = \begin{cases}
    g_0(\vec x) & \text{جے $R_0(\vec{x})$ ہووے} \\
    g_1(\vec x) & \text{جے $R_1(\vec{x})$ ہووے تے $R_0(\vec{x})$ نہ ہووے} \\
    \vdots & \\
    g_{m-1}(\vec x) & \text{جے $R_{m-1}(\vec{x})$ ہووے تے پچھلیاں وچوں کوئی شرط پوری نہ ہووے}
    \\
    g_m(\vec x) & \mbox{ورنہ}
\end{cases}
\]
وی بنیادی بازگشتی اے۔
\end{prop}

\begin{proof}
  جدوں $m = 1$ ہووے، ایہ بس
  ایس طرح متعین فنکشن اے:
  \[
  f(\vec x) = \fn{cond}(\Char{\lnot R_0}(\vec x),g_0(\vec x),g_1(\vec
  x)).
  \]
  جے $m$، $1$ توں وڈا ہووے،
  تاں اِس روپ دیاں تعریفاں
  دی ترکیب کر سکدے آں۔
\end{proof}
\end{document}
